\documentclass[a4paper,10pt]{article}
\usepackage[utf8]{inputenc}
\usepackage{adjustbox}
\usepackage{float}
\usepackage{indentfirst}
\usepackage{listings}
\usepackage{xr}
\usepackage{hyperref}

\externaldocument[CPU-]{cpu}

%opening
\title{
	$SiER^2NA$ \\
	\small Simple Extensible Register-Register Neumann Architecture \\
	\large Instruction Set Architecture \\
	\large 1.0
}
\date{\today}

\begin{document}

\lstset{frame=tb,
	breaklines=true,
	tabsize=4,
	numbers=left,
	basicstyle=\ttfamily
}

\maketitle
\tableofcontents

\section{Preface}

\section{Glossary}
\begin{itemize}
	\item sext18(x) - sign-extend x from bit 17 to 32 bits
	\item sext16(x) - sign-extend x from bit 15 to 32 bits
	\item zext16(x) - zero-extend x from bit 15 to 32 bits
	\item sext8(x) - sign-extend x from bit 7 to 32 bits
	\item zext8(x) - zero-extend x from bit 7 to 32 bits
	\item read8(a) - read 8 bits from address a
	\item read16(a) - read 16 bits from address a
	\item write8(a, x) - write 8 bit value x to address a
	\item write16(a, x) - write 16 bit value x to address a
\end{itemize}


\section{Design}

% TODO
\subsection{Instruction types}
\subsubsection{Register type (R-type)}
\subsubsection{Immediate type (I-type)}
\subsubsection{Immediate jump type (J-type)}

\subsection{Instruction encoding}

\subsubsection{R-type}
\begin{table*}[ht]
	\centering
	\resizebox{\linewidth}{!}{
		\begin{tabular}{|l|l|l|l|l|l|l|l|l|l|l|l|l|l|l|l|l|l|l|l|l|l|l|l|l|l|l|l|l|l|l|l|}
\hline
31 & 30 & 29 & 28 & 27 & 26 & 25 & 24 & 23 & 22 & 21 & 20 & 19 & 18 & 17 & 16 & 15 & 14 & 13 & 12 & 11 & 10 & 9 & 8 & 7 & 6 & 5 & 4 & 3 & 2 & 1 & 0 \\
\hline
\multicolumn{8}{|c|}{ff} & \multicolumn{5}{|c|}{rb} & \multicolumn{5}{|c|}{ra} & \multicolumn{5}{|c|}{ds} & \multicolumn{5}{|c|}{op} & \multicolumn{4}{|c|}{os} \\
\multicolumn{8}{|c|}{function} & \multicolumn{5}{|c|}{register 2} & \multicolumn{5}{|c|}{register 1} & \multicolumn{5}{|c|}{destination} & \multicolumn{5}{|c|}{op code} & \multicolumn{4}{|c|}{op set} \\
\hline
		\end{tabular}
	}
\end{table*}

\subsubsection{I-type}
\begin{table*}[ht]
	\centering
	\resizebox{\linewidth}{!}{
		\begin{tabular}{|l|l|l|l|l|l|l|l|l|l|l|l|l|l|l|l|l|l|l|l|l|l|l|l|l|l|l|l|l|l|l|l|}
\hline
31 & 30 & 29 & 28 & 27 & 26 & 25 & 24 & 23 & 22 & 21 & 20 & 19 & 18 & 17 & 16 & 15 & 14 & 13 & 12 & 11 & 10 & 9 & 8 & 7 & 6 & 5 & 4 & 3 & 2 & 1 & 0 \\
\hline
\multicolumn{2}{|c|}{ff} & \multicolumn{16}{|c|}{vl} & \multicolumn{5}{|c|}{ds} & \multicolumn{5}{|c|}{op} & \multicolumn{4}{|c|}{os} \\
\multicolumn{2}{|c|}{function} & \multicolumn{16}{|c|}{value} & \multicolumn{5}{|c|}{destination} & \multicolumn{5}{|c|}{op code} & \multicolumn{4}{|c|}{op set} \\
\hline
		\end{tabular}
	}
\end{table*}

\subsubsection{J-type}
\begin{table*}[ht]
	\centering
	\resizebox{\linewidth}{!}{
		\begin{tabular}{|l|l|l|l|l|l|l|l|l|l|l|l|l|l|l|l|l|l|l|l|l|l|l|l|l|l|l|l|l|l|l|l|}
\hline
31 & 30 & 29 & 28 & 27 & 26 & 25 & 24 & 23 & 22 & 21 & 20 & 19 & 18 & 17 & 16 & 15 & 14 & 13 & 12 & 11 & 10 & 9 & 8 & 7 & 6 & 5 & 4 & 3 & 2 & 1 & 0 \\
\hline
\multicolumn{5}{|c|}{rc} & \multicolumn{18}{|c|}{tg} & \multicolumn{5}{|c|}{op} & \multicolumn{4}{|c|}{os} \\
\multicolumn{5}{|c|}{comparison result} & \multicolumn{18}{|c|}{target} & \multicolumn{5}{|c|}{op code} & \multicolumn{4}{|c|}{op set} \\
\hline
		\end{tabular}
	}
\end{table*}

\section{Core set (CORE32)}

% os, op, ds, ra, rb, ff, desc
\newcommand{\opr}[7]{
	\begin{table}[H]
		\centering
		\begin{tabular}{|c|c|c|c|c|c|}
	\hline
	os & op & ds & ra & rb & ff \\
	\hline
	#1 & #2 & #3 & #4 & #5 & #6 \\
	\hline
		\end{tabular}
	\end{table}

	#7

	\lstset{title=Implementation}%
}

% os, op, ds, vl, ff, desc
\newcommand{\opi}[6]{
	\begin{table}[H]
		\centering
		\begin{tabular}{|c|c|c|c|c|}
	\hline
	os & op & ds & vl & ff \\
	\hline
	#1 & #2 & #3 & #4 & #5 \\
	\hline
		\end{tabular}
	\end{table}

	#6

	\lstset{title=Implementation}%
}

% os, op, tg, rc, desc
\newcommand{\opj}[5]{
	\begin{table}[H]
		\centering
		\begin{tabular}{|c|c|c|c|}
	\hline
	os & op & tg & rc \\
	\hline
	#1 & #2 & #3 & #4 \\
	\hline
		\end{tabular}
	\end{table}

	#5

	\lstset{title=Implementation}%
}

\subsection{Miscellaneous}
\subsubsection{nop - no operation}
\opr{0x00}{0x00}{--}{--}{--}{--}
	{Performs no operation. The number of cycles consumed by this instruction is implementation-dependent and not guaranteed.}

\subsubsection{cpuid - CPU identification}
\opr{0x00}{0x01}{result}{--}{--}{--}
	{Writes the CPUID structure into \emph{result}. See \hyperref[CPU-st:CPUID]{CPU: CPUID}}

\begin{lstlisting}
// implementation-specific
\end{lstlisting}

\subsection{Load/Store operations}
\subsubsection{ldw - load word}
\opr{0x00}{0x02}{destination}{\&source}{--}{0x00}
	{Loads a single word from memory at \emph{source} into \emph{destination}}
\begin{lstlisting}
value = *source
\end{lstlisting}

\subsubsection{ldhz - load half-word zero-extended}
\opr{0x00}{0x02}{destination}{\&source}{--}{0x01}
	{Loads a zero-extended half-word from memory at \emph{source} into \emph{destination}}
\begin{lstlisting}
value = zext16(read16(*source))
\end{lstlisting}

\subsubsection{ldhs - load half-word sign-extended}
\opr{0x00}{0x02}{destination}{\&source}{--}{0x02}
	{Loads a sign-extended half-word from memory at \emph{source} into \emph{destination}}
\begin{lstlisting}
value = sext16(read16(*source))
\end{lstlisting}

\subsubsection{ldbz - load byte zero-extended}
\opr{0x00}{0x02}{destination}{\&source}{--}{0x03}
	{Loads a zero-extended byte from memory at \emph{source} into \emph{destination}}
\begin{lstlisting}
value = zext8(read8(*source))
\end{lstlisting}

\subsubsection{ldbs - load byte sign-extended}
\opr{0x00}{0x02}{destination}{\&source}{--}{0x04}
	{Loads a sign-extended byte from memory at \emph{source} into \emph{destination}}
\begin{lstlisting}
value = sext8(read8(*source))
\end{lstlisting}

\subsubsection{lvi - load sign-extended value immediate}
\opi{0x00}{0x03}{destination}{value}{--}
	{Loads \emph{value} into the lower half-word of \emph{destination}, sign-extending it to the upper-half}
\begin{lstlisting}
destination = sext16(value)
\end{lstlisting}

\subsubsection{lli - load lower half-word immediate}
\opi{0x00}{0x04}{destination}{value}{0x01}
	{Loads \emph{value} into the lower half-word of \emph{destination}}
\begin{lstlisting}
destination = (*destination & 0xFFFF0000) | value
\end{lstlisting}

\subsubsection{lui - load upper half-word immediate}
\opi{0x00}{0x04}{destination}{value}{0x02}
	{Loads \emph{value} into the upper half-word of \emph{destination}}
\begin{lstlisting}
destination = (*destination & 0x0000FFFF) | (value << 16)
\end{lstlisting}

\subsubsection{stw - store word}
\opr{0x00}{0x05}{\&destination}{source}{--}{0x00}
	{Writes a single word into memory at \emph{destination} from \emph{source}}
\begin{lstlisting}
*destination = source
\end{lstlisting}

\subsubsection{sth - store half-word}
\opr{0x00}{0x05}{\&destination}{source}{--}{0x01}
	{Writes a half-word into memory at \emph{destination} from \emph{source}}
\begin{lstlisting}
*destination = source >> 16
\end{lstlisting}

\subsubsection{stb - store byte}
\opr{0x00}{0x05}{\&destination}{source}{--}{0x02}
	{Writes a byte into memory at \emph{destination} from \emph{source}}
\begin{lstlisting}
*destination = source >> 24
\end{lstlisting}

\subsection{Register operations}
\subsubsection{mov - move}
\opr{0x00}{0x06}{destination}{source}{--}{--}
	{Copies the value of \emph{source} into \emph{destination}}
\begin{lstlisting}
destination = source
\end{lstlisting}

\subsection{Control flow}
\subsubsection{ja - jump absolute}
\opr{0x00}{0x07}{\&target}{--}{--}{--}
	{Sets the \emph{instruction pointer} to the value of \emph{target}}
\begin{lstlisting}
rip = target
\end{lstlisting}

\subsubsection{jae - jump absolute if equal}
\opr{0x00}{0x08}{\&target}{cmp}{--}{0x00}
	{Sets the \emph{instruction pointer} to the value of \emph{target} if the comparison result \emph{cmp} is \textbf{EQUAL}}
\begin{lstlisting}
if (cmp & 0b0001) == 1 {
	rip = target
}
\end{lstlisting}

\subsubsection{jane - jump absolute if not equal}
\opr{0x00}{0x08}{\&target}{cmp}{--}{0x01}
	{Sets the \emph{instruction pointer} to the value of \emph{target} if the comparison result \emph{cmp} is \textbf{NOT EQUAL}}
\begin{lstlisting}
if (cmp & 0b0010) == 1 {
	rip = target
}
\end{lstlisting}

\subsubsection{jalt - jump absolute if less than}
\opr{0x00}{0x08}{\&target}{cmp}{--}{0x02}
	{Sets the \emph{instruction pointer} to the value of \emph{target} if the comparison result \emph{cmp} is \textbf{LESS THAN}}
\begin{lstlisting}
if (cmp & 0b0100) == 1 {
	rip = target
}
\end{lstlisting}

\subsubsection{jage - jump absolute if greater or equal}
\opr{0x00}{0x08}{\&target}{cmp}{--}{0x03}
	{Sets the \emph{instruction pointer} to the value of \emph{target} if the comparison result \emph{cmp} is \textbf{GREATER THAN} or \textbf{EQUAL}}
\begin{lstlisting}
if (cmp & 0b1000) == 1 || (cmp & 0b0001) == 1 {
	rip = target
}
\end{lstlisting}

\subsubsection{jr - jump relative}
\opj{0x00}{0x09}{target}{--}
	{Sets the \emph{instruction pointer} to sign-extended \emph{target} relative to the current value of the \emph{instruction pointer}}
\begin{lstlisting}
rip = rip + sext18(target)
\end{lstlisting}

\subsubsection{jre - jump relative if equal}
\opj{0x00}{0x0A}{target}{cmp}
	{Sets the \emph{instruction pointer} to sign-extended \emph{target} relative to the current value of the \emph{instruction pointer} if the comparison result \emph{cmp} is \textbf{EQUAL}}
\begin{lstlisting}
if (cmp & 0b0001) == 1 {
	rip = rip + sext18(target)
}
\end{lstlisting}

\subsubsection{jrne - jump relative if not equal}
\opj{0x00}{0x0B}{target}{cmp}
	{Sets the \emph{instruction pointer} to sign-extended \emph{target} relative to the current value of the \emph{instruction pointer} if the comparison result \emph{cmp} is \textbf{NOT EQUAL}}
\begin{lstlisting}
if (cmp & 0b0010) == 1 {
	rip = rip + sext18(target)
}
\end{lstlisting}

\subsubsection{jrlt - jump relative if less than}
\opj{0x00}{0x0C}{target}{cmp}
	{Sets the \emph{instruction pointer} to sign-extended \emph{target} relative to the current value of the \emph{instruction pointer} if the comparison result \emph{cmp} is \textbf{LESS THAN}}
\begin{lstlisting}
if (cmp & 0b0100) == 1 {
	rip = rip + sext18(target)
}
\end{lstlisting}

\subsubsection{jrge - jump relative if greater or equal}
\opj{0x00}{0x0D}{target}{cmp}
	{Sets the \emph{instruction pointer} to sign-extended \emph{target} relative to the current value of the \emph{instruction pointer} if the comparison result \emph{cmp} is \textbf{GREATER THAN} or \textbf{EQUAL}}
\begin{lstlisting}
if (cmp & 0b1000) == 1 || (cmp & 0b0001) == 1 {
	rip = rip + sext18(target)
}
\end{lstlisting}

\subsubsection{calla - call procedure absolute}
\opr{0x00}{0x0E}{\&target}{--}{--}{--}
	{Sets the \emph{procedure tail pointer} to the value of the \emph{instruction pointer} and then sets the \emph{instruction pointer} to the value of \emph{target}}
\begin{lstlisting}
rtp = rip + 4
rip = target
\end{lstlisting}

\subsubsection{callr - call procedure relative immediate}
\opj{0x00}{0x0F}{\&target}{--}
	{Sets the \emph{procedure tail pointer} to the value of the \emph{instruction pointer} and then sets the \emph{instruction pointer} to the value of sign-extended \emph{target} relative to the current value of the \emph{instruction pointer}}
\begin{lstlisting}
rtp = rip + 4
rip = rip + sext18(target)
\end{lstlisting}

\subsubsection{ret - return from procedure}
\opr{0x00}{0x10}{--}{--}{--}{--}
	{Sets the \emph{instruction pointer} to the value of the \emph{procedure tail pointer}}
\begin{lstlisting}
rip = rtp
\end{lstlisting}

\subsection{Mathematical operations}
\subsubsection{cmp - compare values}
\opr{0x00}{0x11}{result}{a}{b}{--}
	{
		Compares the values of \emph{a} and \emph{b} and stores the comparison result in \emph{result}.

		The comparison result is a 4-bit flag value and maps to the following results:
		\begin{table}[H]
			\centering
			\begin{tabular}{|c|c|c|}
				\hline
				bit & mnemonic & description \\
				\hline
				0 & EQUAL & if set, \emph{a} and \emph{b} are equal \\
				\hline
				1 & NOT EQUAL & if set, \emph{a} and \emph{b} are not equal \\
				\hline
				2 & LESS THAN & if set, \emph{a} is less than \emph{b} \\
				\hline
				3 & GREATER THAN & if set, \emph{a} is greater than \emph{b} \\
				\hline
			\end{tabular}
		\end{table}
	}
\begin{lstlisting}
result = 0

if a == b {
	result = result | (1 << 0)
} else {
	result = result | (1 << 1)

	if a < b {
		result = result | (1 << 2)
	}

	if a > b {
		result = result | (1 << 3)
	}
}
\end{lstlisting}

\subsubsection{add - integer addition}
\opr{0x00}{0x12}{result}{a}{b}{0x00}
	{Adds the values of \emph{a} and \emph{b} together and stores the result in \emph{result}}
\begin{lstlisting}
result = a + b
\end{lstlisting}

\subsubsection{sub - integer subtraction}
\opr{0x00}{0x12}{result}{a}{b}{0x01}
	{Subtracts the value of \emph{b} from the value of \emph{a} and stores the result in \emph{result}}
\begin{lstlisting}
result = a - b
\end{lstlisting}

\subsubsection{addi - integer addition sign-extended immediate}
\opi{0x00}{0x13}{a}{b}{--}
	{Adds the value of sign-extended \emph{b} and the value of \emph{a} together, storing the result in \emph{a}}
\begin{lstlisting}
a = a + sext16(b)
\end{lstlisting}

\subsubsection{mul - integer multiplication}
\opr{0x00}{0x14}{result}{a}{b}{0x00}
	{Multiplies the values of \emph{a} and \emph{b} together and stores the result in \emph{result}}
\begin{lstlisting}
result = a * b
\end{lstlisting}

\subsubsection{div - integer division}
\opr{0x00}{0x14}{result}{a}{b}{0x01}
	{Divides the value of \emph{a} by the value of \emph{b} stores the result in \emph{result}
		\begin{itemize}
			\item Fractional results are truncated towards 0
			\item Division by 0 raises \textbf{\#DBZ}
		\end{itemize}
	}
\begin{lstlisting}
result = a // b
\end{lstlisting}

\subsubsection{mod - integer remainder}
\opr{0x00}{0x14}{result}{a}{b}{0x02}
	{Stores the division remainder of dividing \emph{a} by \emph{b} in \emph{result}
		\begin{itemize}
			\item Fractional results are truncated towards 0
			\item Modulo by 0 throws a DIVISION BY ZERO exception
		\end{itemize}
	}
\begin{lstlisting}
result = a % b
\end{lstlisting}

\subsection{Bitwise operations}
\subsubsection{and - bitwise AND}
\opr{0x00}{0x15}{result}{a}{b}{0x00}
	{Performs a bitwise AND operation on the values of \emph{a} and \emph{b} and stores the result in \emph{result}}
\begin{lstlisting}
result = a & b
\end{lstlisting}

\subsubsection{or - bitwise OR}
\opr{0x00}{0x15}{result}{a}{b}{0x01}
	{Performs a bitwise OR operation on the values of \emph{a} and \emph{b} and stores the result in \emph{result}}
\begin{lstlisting}
result = a | b
\end{lstlisting}

\subsubsection{xor - bitwise XOR}
\opr{0x00}{0x15}{result}{a}{b}{0x02}
	{Performs a bitwise XOR operation on the values of \emph{a} and \emph{b} and stores the result in \emph{result}}
\begin{lstlisting}
result = a ^ b
\end{lstlisting}

\subsubsection{shl - logical shift left}
\opr{0x00}{0x16}{result}{a}{n}{0x00}
	{Shifts the value of \emph{a} left by \emph{n} places and stores the result in \emph{result}}
\begin{lstlisting}
result = a << n
\end{lstlisting}

\subsubsection{shr - logical shift right}
\opr{0x00}{0x16}{result}{a}{n}{0x01}
	{Shifts the value of \emph{a} right by \emph{n} places and stores the result in \emph{result}}
\begin{lstlisting}
result = a >> n
\end{lstlisting}

\subsubsection{shsr - arithmetic (signed) shift right}
\opr{0x00}{0x16}{result}{a}{n}{0x02}
	{Shifts the value of \emph{a} right by \emph{n} places and stores the result in \emph{result}, sign-replicating it}
\begin{lstlisting}
result = (a >> n) | (a >> n >> (31 - n) << 31)
\end{lstlisting}

\subsubsection{shli - logical shift left immediate}
\opi{0x00}{0x17}{a}{n}{0x00}
	{Shifts the value of \emph{a} in-place left by \emph{n} places}
\begin{lstlisting}
a = a << zext16(n)
\end{lstlisting}

\subsubsection{shri - logical shift right immediate}
\opi{0x00}{0x17}{a}{n}{0x01}
	{Shifts the value of \emph{a} in-place right by \emph{n} places}
\begin{lstlisting}
a = a >> zext16(n)
\end{lstlisting}

\subsubsection{shsri - arithmetic (signed) shift right immediate}
\opi{0x00}{0x17}{a}{n}{0x02}
	{Shifts the value of \emph{a} in-place right by \emph{n} places, sign-replicating the result}
\begin{lstlisting}
a = a >> (a >> zext16(n)) | (a >> zext16(n) >> (31 - zext16(n)) << 31)
\end{lstlisting}

\subsection{Interrupt control}
\subsubsection{int - raise software interrupt}
\opr{0x00}{0x18}{id}{--}{--}{--}
	{Raises a software interrupt}

\subsubsection{inti - raise software interrupt immediate}
\opi{0x00}{0x19}{id}{--}{--}
	{Raises a software interrupt}

\subsubsection{intm - mask interrupt}
\opr{0x00}{0x1A}{id}{enable}{--}{--}
	{If \emph{enable} is non-zero masks an interrupt, telling the CPU not to handle it through a \hyperref[CPU-ir:HANDLER]{Interrupt Routine Handler}. Otherwise, (re-)enables it.}
\begin{lstlisting}
if enable {
	rrap[id * 0xC] = rrap[id * 0xC] | 1
} else {
	rrap[id * 0xC] = rrap[id * 0xC] & ~1
}
\end{lstlisting}
\end{document}
